% Clase 15 — resistores en serie y en paralelo. Mismos R1 y R2; cambia la
% conexión. En serie pasa la misma corriente y se suman las caídas; en paralelo
% la caída es la misma y se suman las corrientes (ley de los nudos).
% Las flechas de corriente se dibujan aparte, no con las etiquetas `i=` de
% circuitikz.
\begin{tikzpicture}[scale=1]

  % =============== (a) serie ===============
  \begin{scope}[line width=0.9pt, color=figblue]
    \draw (0,1) -- (0.6,1);
    \draw (0.6,1) to[R] (2.6,1);
    \draw (2.6,1) to[R] (4.6,1);
    \draw (4.6,1) -- (5.2,1);
    \draw[fill=white] (0,1) circle (1.6pt);
    \draw[fill=white] (5.2,1) circle (1.6pt);
  \end{scope}
  \node[figlbl, figblue] at (1.6,1.55) {$R_1$};
  \node[figlbl, figblue] at (3.6,1.55) {$R_2$};
  \draw[figvecres] (0.05,0.55) -- (0.75,0.55);
  \node[figlbl, text=figred] at (0.4,0.22) {$I$};
  \draw[figvecres] (4.5,0.55) -- (5.2,0.55);
  \node[figlbl, text=figred] at (4.85,0.22) {$I$};
  \node[figlbl, anchor=north] at (2.6,-0.75) {(a) serie};

  % =============== (b) paralelo ===============
  \begin{scope}[xshift=7.2cm]
    \begin{scope}[line width=0.9pt, color=figblue]
      \draw (0,1) -- (0.9,1);
      \draw (0.9,1) -- (0.9,1.9);
      \draw (0.9,1.9) to[R] (3.5,1.9);
      \draw (3.5,1.9) -- (3.5,1);
      \draw (0.9,1) -- (0.9,0.1);
      \draw (0.9,0.1) to[R] (3.5,0.1);
      \draw (3.5,0.1) -- (3.5,1);
      \draw (3.5,1) -- (4.4,1);
      \fill (0.9,1) circle (2pt);
      \fill (3.5,1) circle (2pt);
      \draw[fill=white] (0,1) circle (1.6pt);
      \draw[fill=white] (4.4,1) circle (1.6pt);
    \end{scope}
    \node[figlbl, figblue] at (2.2,2.45) {$R_1$};
    \node[figlbl, figblue] at (2.2,-0.45) {$R_2$};
    \draw[figvecres] (0.05,1.3) -- (0.7,1.3);
    \node[figlbl, text=figred] at (0.35,1.62) {$I$};
    \draw[figvecres] (1.0,2.2) -- (1.6,2.2);
    \node[figlbl, text=figred, anchor=east] at (0.9,2.2) {$I_1$};
    \draw[figvecres] (1.0,-0.2) -- (1.6,-0.2);
    \node[figlbl, text=figred, anchor=east] at (0.9,-0.2) {$I_2$};
    \node[figlbl, anchor=north] at (2.2,-0.75) {(b) paralelo};
  \end{scope}

\end{tikzpicture}
